\begin{pastexam}[label={px:ntu-114-dsa-05}]{臺大資網所 114 學年度 第 5 題}
（單選；答對 5 分，不答或答錯 0 分）

(5 points) Many graph theory can be used to find the number of a chemical structure; a possible one can be used
according to the following steps:\\
a) Represent the molecule as an undirected graph $G=(V,E)$, where $V$ is the set of vertices and $E$ is the set of
edges (bonds).\\
b) Initialize each vertex $v\in V$ with a value equal to its degree (number of incident edges).\\
c) For a fixed number of iterations or until convergence:\\
\hspace*{1em}a. For each vertex $v$, compute a new value as the sum of the current values of its neighbors (not
including itself).\\
\hspace*{1em}b. Update all vertex values simultaneously.\\
d) After the iterations, sort the vertices based on their final values in non-increasing order, tie breaking with
smaller labels.

Consider the following 2-Carboxyoxybenzoic acid structure represented as a graph (Numbers in parentheses are
temporary labels for reference, not the final numbering.):
\begin{center}
\begin{tikzpicture}[x=1.05cm,y=1.05cm,font=\small,line cap=round,
  bond/.style={line width=0.7pt,shorten >=6pt,shorten <=6pt},
  bondb/.style={line width=0.7pt,shorten >=8.5pt,shorten <=8.5pt}]
  % 原子座標（依原圖位置）
  \coordinate (O12) at (1.50,-0.85); \coordinate (O13) at (3.35,-0.85);
  \coordinate (C11) at (2.45,-1.50); \coordinate (C10) at (2.45,-2.50);
  \coordinate (C9)  at (1.40,-3.10); \coordinate (C8)  at (1.40,-4.15);
  \coordinate (C7)  at (2.45,-4.80); \coordinate (C6)  at (3.55,-4.15);
  \coordinate (C5)  at (3.55,-3.10); \coordinate (O4)  at (4.55,-2.50);
  \coordinate (C3)  at (5.65,-3.10); \coordinate (O1)  at (6.60,-2.50);
  \coordinate (O2)  at (5.65,-4.15);
  % 單鍵
  \foreach \a/\b in {C11/O13,C11/C10,C9/C8,C7/C6,C5/C10,C5/O4,O4/C3,C3/O1}
    \draw[bond] (\a)--(\b);
  % 雙鍵（主線 + 平行短線）
  \foreach \a/\b in {O12/C11,C10/C9,C8/C7,C5/C6,C3/O2}{
    \draw[bond] (\a)--(\b);
    \draw[bondb] ($(\a)!2.6pt!90:(\b)$)--($(\b)!2.6pt!-90:(\a)$);
  }
  % 原子與暫時編號
  \foreach \p/\s/\n in {O12/O/12,O13/O/13,C11/C/11,C10/C/10,C9/C/9,C8/C/8,C7/C/7,C6/C/6,C5/C/5,O4/O/4,C3/C/3,O1/O/1,O2/O/2}{
    \node[inner sep=0.5pt] at (\p) {\s};
    \node[anchor=west,inner sep=0pt,font=\footnotesize] at ($(\p)+(0.17,0)$) {(\n)};
  }
\end{tikzpicture}
\end{center}
Apply the above algorithm for two iterations to this graph. What is the correct ordering of the vertices after
these iterations?
\begin{enumerate}[label=(\Alph*)]
\item 5, 6, 10, 11, 9, 4, 3, 7, 8, 12, 13, 1, 2
\item 5, 10, 11, 9, 3, 6, 4, 7, 8, 12, 13, 1, 2
\item 5, 6, 10, 11, 9, 3, 4, 7, 8, 12, 13, 1, 2
\item 5, 10, 11, 9, 6, 3, 4, 7, 8, 12, 13, 1, 2
\item None of the above choices.
\end{enumerate}
\end{pastexam}

\begin{pastsoln}
答案是 (E)。 這是化學上替分子編號的 Morgan 演算法。

\textbf{建圖}。 題目說 $E$ 是鍵的集合，雙鍵也只是一條邊（圖是簡單圖），度數就是相連的原子數。 從圖讀出 $13$ 條邊：
\[
\begin{aligned}
&12\text{--}11,\ 11\text{--}13,\ 11\text{--}10,\ 10\text{--}9,\ 10\text{--}5,\ 9\text{--}8,\ 8\text{--}7,\\
&7\text{--}6,\ 6\text{--}5,\ 5\text{--}4,\ 4\text{--}3,\ 3\text{--}1,\ 3\text{--}2 .
\end{aligned}
\]
$5, 6, 7, 8, 9, 10$ 是苯環。

\textbf{迭代}。 第 $0$ 輪是度數；之後每一輪，每個頂點的新值是\textbf{上一輪}所有鄰居的值相加（同時更新）。
\begin{center}
\small
\begin{tabular}{c|ccccccccccccc}
頂點 & 1 & 2 & 3 & 4 & 5 & 6 & 7 & 8 & 9 & 10 & 11 & 12 & 13 \\ \hline
鄰居 & 3 & 3 & 1,2,4 & 3,5 & 4,6,10 & 5,7 & 6,8 & 7,9 & 8,10 & 5,9,11 & 10,12,13 & 11 & 11 \\ \hline
第 0 輪 & 1 & 1 & 3 & 2 & 3 & 2 & 2 & 2 & 2 & 3 & 3 & 1 & 1 \\
第 1 輪 & 3 & 3 & 4 & 6 & 7 & 5 & 4 & 4 & 5 & 8 & 5 & 3 & 3 \\
第 2 輪 & 4 & 4 & 12 & 11 & 19 & 11 & 9 & 9 & 12 & 17 & 14 & 5 & 5
\end{tabular}
\end{center}
幾個例子：第 1 輪的 $10$ 是 $5, 9, 11$ 第 0 輪的值 $3 + 2 + 3 = 8$；第 2 輪的 $5$ 是 $4, 6, 10$ 第 1 輪的值
$6 + 5 + 8 = 19$，$3$ 是 $1, 2, 4$ 的 $3 + 3 + 6 = 12$，$9$ 是 $8, 10$ 的 $4 + 8 = 12$。

\textbf{排序}。 依第 2 輪的值由大到小，同值時標號小的在前：
\[
5\,(19),\ 10\,(17),\ 11\,(14),\ 3\,(12),\ 9\,(12),\ 4\,(11),\ 6\,(11),\ 7\,(9),\ 8\,(9),\ 12\,(5),\ 13\,(5),\ 1\,(4),\ 2\,(4).
\]
和選項比：
\begin{itemize}
\item (A)、(C) 第二個是 $6$，但 $6$ 只有 $11$，排在 $10$、$11$ 之後。
\item (B) 是 $5, 10, 11, 9, 3, \ldots$：$9$ 和 $3$ 都是 $12$，同值時標號小的在前，應該是 $3$ 在 $9$ 前；後面 $6, 4$
也一樣，應該是 $4, 6$。
\item (D) 把 $6$ 排在 $3$ 前面，但 $6$ 是 $11$、$3$ 是 $12$。
\end{itemize}
四個都不對，選 (E)。 如果把雙鍵算成兩條邊（多重圖），第 2 輪的值會變成 $5{:}\,44$、$10{:}\,43$、$9{:}\,37$、
$11{:}\,34$、\ldots，順序也不是 (A)--(D) 的任何一個，答案仍是 (E)。
\end{pastsoln}
